\section{Economic Simulation}

When attempting to simulate a consumer market, there are quite a few approaches
you could choose between. The vast majority of models fall into 2 categories:

\begin{enumerate}
	\item{Statistical models: models where the simulation uses regressions and
	      other statistical analysis tools to determine outcomes.}
	\item{Agent-based models: these models choose to, instead of simulating the
	      whole market at once, instantiate individual \textit{agents} that have
	      the ability to independently make decisions.  }
\end{enumerate}

There are 2 main advantages to statistical models: they are simpler to build, as
you can often base them around simple regressions while still returning convincing
results, and performance, unless the implementer makes grave mistakes, statistical
models almost always outperform (in terms of speed) agent-based models of the same
caliber.

For example, the cult classic city builder SimCity 4 was able to simulate regions
with populations up to 100 million with very minor performance issues,
\footnote{
	This Kotaku article from 2014 shows off a city of over 100 million
	inhabitants, the Guinness world record, without any performance problems.

	\noindent  Link to article:
	\href{https://kotaku.com/simcity-megacity-has-over-100-000-000-million-people-li-1629131314}
	{https://kotaku.com/simcity-megacity-has-over-100-000-000-million-people-li-1629131314};

	\noindent World record:
	\href{https://www.guinnessworldrecords.com/world-records/418977-largest-simcity-4-city-region}
	{https://www.guinnessworldrecords.com/world-records/418977-largest-simcity-4-city-region}
}
while over 20 years later, games like Cities: Skylines 2 are still plagued by
performance issues from simulating just 100'000 agents on modern multi-core
hardware.

The biggest downside of statistical modelling for this project is flexibility.
An agent-based model can simulate virtually anything, should the designer choose
to, without necessarily major changes. This was the deciding factor for me, as
the development of the game would take time, during which I would inevitably want
to rework aspects of the simulation to optimise certain aspects.

\subsection{The Market}

\subsubsection*{The customers}

The customer is simulated as a bundle of different characteristics, preferences,
and needs; these all combine to produce unique behaviors for each individual
agent. Taking inspiration from Economy Week, each customer has the following
properties: \footnote{The names provided the table may not be accurate to the
	corresponding implementations found in the code to simplify the explanation,
	as the implementation contains subtle differences for performance and
	technological reasons.
}

{\small
	\begin{tabular}{l l p{0.45\textwidth}} \toprule
		Property             & Type               & Description                        \\ \midrule

		Base Need            & Number             & The amount of products a customer
		wants. When their existing products wear out, they buy more.                   \\

		Owned Products       & List[number]       & Keeps track of the remaining
		durability of owned products.                                                  \\

		Purchasing Threshold & Number             & How attractive a product has to
		be for the customer to buy it.                                                 \\

		Savviness            & Number             & How well informed the customer is. \\

		Loyalties            & List[number]       & How loyal the customer is to a
		certain company.                                                               \\

		Loyalty Factor       & Number             & How much the customer's Loyalties
		influence their                                                                \\

		Preferences          & Dictionary[number] & Which aspects of a product impact
		the customer's decisions the most. For example, a customer could value low
		prices more than they do the quality of the product.                           \\ \bottomrule
	\end{tabular}}

Originally the preferences were determined by sampling a \textit{weighted normal
	distribution}; however, after some playtesting and reading
\href{https://en.wikipedia.org/wiki/The_Tipping_Point}{ The Tipping Point } by
Malcom Gladwell, I reworked the generation of customers by instead creating
\textit{archetypes} and slightly modifying their preferences randomly.

Each of the archetypes has some preferences boosted and others diminished. For
example, "the enthusiast" will not shy away from buying an expensive product,
even multiple, if they perceive it to be unique and high quality. On the other
hand, "the Everyman" will prefer a simpler, cheaper, long-lasting product.

	{\small archetypes have not yet been fleshed out}

\subsubsection*{The Companies}

During the decisions phase, companies were able to develop products and set up
production chains. Now the products made by companies are put on the market
along with corresponding marketing campaigns.

\subsubsection*{External Factors}

Since the simulation is limited to the customers and products, external factors
are able to add extra challenge and variety to the game. There are external
factors that influence, for example, how easily people will spend money,
resembling the \href{https://en.wikipedia.org/wiki/Credit_cycle}
{ credit cycle } experienced in real life, where after a period of high
consumer spending there is a period of low spending while consumers save back up
and pay off loans.

\subsection{The Marketing Phase}

In this phase, the marketing campaigns of companies are simulated. Each
marketing campaign has a style, a quality and a quantity.

The style is determined by the player, which can influence the perception of the
product. This can be used to round out the \textit{less} attractive parts of the
product. For example, if a player had a high quality product that they wanted to
market to a more eco-focused audience, instead of having to undertake the
difficult and expensive process of actually producing in an environmentally
friendly manner, they can just focus on promoting the product as eco-friendly.
\footnote{This technique is often labelled as
	\href{https://en.wikipedia.org/wiki/Greenwashing}{Greenwashing}.
	The most famous example is Volkswagen's "clean diesel" emissions scandal in 2015
	where Volkswagen programmed diesel engined cars to run efficiently when they
	were being tested while reverting to "dirty" operation when driven on the road.

	\noindent Further reading:
	\href{https://en.wikipedia.org/wiki/Volkswagen_emissions_scandal}
	{https://en.wikipedia.org/wiki/Volkswagen\_emissions\_scandal}
}

The quality of a campaign influences how much the perception of the product is
changed. This is mostly determined by the skill and motivation of the companies
marketing staff.

Lastly, the quantity of a campaign, expressed in a monetary amount, determines
who even sees the campaign. A high quantity of marketing means there is a high
chance that people know about a product. A campaign could theoretically be
extremely effective, however if no-one sees it, it isn’t actually successful
since there are still few people who even know about the product’s existence.

All of these factors are mixed in with each customer's "savviness" stat. People
with more savviness are generally more informed on a market and thus have a
higher chance of knowing about a product, however they are also less effected
by marketing and thus perceive it more objectively. Other people with lower
savviness will easily be swayed by effective marketing.

\subsection{The Purchasing Phase}

After the effects of marketing are evaluated, if the customer needs a product,
it will evaluate the dot product of its preferences and the product's properties
(the decision factors) for each product.

$$
	\text{decision factors} = \begin{pmatrix}
		\text{ quality } \\
		...              \\
		\text{ durability }
	\end{pmatrix} \cdot \begin{pmatrix}
		\text{ quality preference } \\
		...                         \\
		\text{ durability preference }
	\end{pmatrix}
$$

The customer then checks if the product is still available; if so, they try to
buy enough to fulfill their needs.
